% PraTeXの横組日本語classの代表入力。
\documentclass[a4paper,10pt]{prjsarticle}

% classは本文・表題にupjisr-h、見出しにupjisg-hを選択する。
% 別JFMを試す時だけ、次の任意adapterがhookを後から置き換える。
% 別engineのversion primitiveを定義して分岐を偽装しない。
\InputIfFileExists{prjsarticle-test-adapter.tex}{}{}

\title{PraTeXによる日本語とLatin textの混植}
\author{PraTeX開発チーム}
\date{2026年8月23日}

\begin{document}
\maketitle

\begin{abstract}
この文書は、横組の日本語段落、欧文、見出し、基本的な箇条書きを一つの
classで確認するための標本である。PraTeX and Latin text share the same line.
\end{abstract}

\section{和欧混植の段落}
日本語の文中にPraTeX、W3C JLReq、safe Rustのような欧文を置く。
句読点を含む行分割と和欧文間の自動空白は、classのmacroではなく
PraTeX engine coreが決定する。

段落頭は一字下げとし、段落間には追加の空きを置かない。二つ目の段落でも
日本語とLatin wordsを自然に混在させる。

\subsection{基本的な箇条書き}
\begin{itemize}
  \item 日本語の項目とLatin textを同じ行に置く。
  \item 項目間隔は本文の行送りを基準にする。
\end{itemize}

\begin{enumerate}
  \item JFM/TFMから和文glyphを作る。
  \item K/X spacingと禁則を中央finalizerで適用する。
\end{enumerate}

\begin{description}
  \item[PraTeX] 日本語・欧文・和欧混植を一級機能として扱うengine。
  \item[prjsarticle] PraTeX identityだけを要求する横組class。
\end{description}

\end{document}
